\documentclass[10pt]{article}
\usepackage[utf8]{inputenc}
\usepackage{amsmath, amssymb}
\usepackage{graphicx}
\usepackage{hyperref}
\usepackage{geometry}
\geometry{a4paper, margin=1in}

\title{EC527 Final Project Report: GMM Optimization}
\author{Author Name(s)}
\date{\today}

\begin{document}

\maketitle

\tableofcontents
\newpage

\section{Description of the Problem}
This project focuses on the performance optimization of the Gaussian Mixture Model (GMM) training process. GMM is a robust probabilistic model used extensively in machine learning, signal processing, and data analysis for tasks such as clustering, anomaly detection, and density estimation. Training a GMM entails fitting a dataset of $N$ points in a $D$-dimensional space to $K$ distinct Gaussian distributions by estimating their corresponding weights, means, and covariance matrices.

The parameter estimation is heavily reliant on the iterative Expectation-Maximization (EM) algorithm, which alternates between two computationally rigorous phases:
\begin{itemize}
    \item \textbf{E-Step (Expectation):} Evaluates the multivariate Gaussian probability density function for every data point across all $K$ clusters to compute posterior probabilities (responsibilities).
    \item \textbf{M-Step (Maximization):} Aggregates the responsibilities across the dataset to update the model parameters (weights, means, and covariances).
\end{itemize}

As dataset sizes ($N$) and dimensionalities ($D$) increase, the EM algorithm becomes highly computationally and memory intensive due to continuous data passes and complex linear algebra operations (such as Cholesky decompositions for covariance matrix inversion). The primary objective of this project is to analyze the performance bottlenecks of a serial GMM implementation and accelerate the algorithm by utilizing parallel computing paradigms, specifically OpenMP for multi-core CPUs and CUDA for GPU architectures.

\section{Serial Code and Algorithm}
The base serial implementation fundamentally loops through a predefined number of iterations or until convergence, executing the Expectation (E) and Maximization (M) steps sequentially. 

\begin{itemize}
    \item \textbf{Expectation Step:} For every data point $x_i$ ($1 \le i \le N$) and every mixture component $k$ ($1 \le k \le K$), the algorithm computes the continuous multidimensional Gaussian probability. This involves first performing a Cholesky decomposition to invert the covariance matrix, then computing the Mahalanobis distance by subtracting the mean vector, multiplying by the inverse covariance matrix, and calculating the exponential. The results are normalized to determine the responsibility $\gamma_{i,k}$, representing the probabilistic assignment of point $x_i$ to cluster $k$.
    \item \textbf{Maximization Step:} The responsibilities are used to recalculate the parameters of the distributions. First, the algorithm computes the total effective weight for each cluster. Next, it performs a full sweep through the dataset to compute the new mean vectors. Finally, it executes another sweep to update the covariance matrices using the newly found means and the responsibility weights.
\end{itemize}

\textbf{Algorithmic Complexity:} The computational complexity per EM iteration is driven by varying operations depending on the dataset grid. During the E-step, computing the multivariate Gaussian PDF requires a Cholesky decomposition and subsequent matrix inversion of the covariance matrix for each of the $K$ clusters, an operation that scales dramatically as $\mathcal{O}(K \cdot D^3)$. Furthermore, computing the Mahalanobis distance for all $N$ points and updating the covariances in the M-step dictates vector-matrix multiplications scaling as $\mathcal{O}(N \cdot K \cdot D^2)$.

Therefore, for small dimensionality ($D \le 16$), the $\mathcal{O}(N \cdot K \cdot D^2)$ memory-bound operations dominate the execution. As $D$ increases, the $\mathcal{O}(K \cdot D^3)$ Cholesky matrix inversions cause the workload to become immensely compute-bound, representing the hardware inflection point between CPU caching/vectorization advantages and raw GPU FLOP throughput. Over $I$ iterations, the total computational complexity is bounded by $\mathcal{O}(I \cdot (K \cdot D^3 + N \cdot K \cdot D^2))$.

\section{Performance Profiling \& Bottlenecks}
Performance profiling on the serial execution (using tools like \texttt{gprof} or Xcode Instruments) reveals that the overwhelming majority of the runtime is concentrated within two heavily nested loops: the Gaussian probability calculation in the E-step, and the covariance matrix calculation in the M-step. 

When evaluating the arithmetic intensity (the ratio of floating-point operations to bytes of memory accessed), GMM training is notoriously \textbf{memory-bandwidth bound} rather than strictly compute-bound. 

For every iteration:
\begin{enumerate}
    \item The E-step reads the entire $N \times D$ dataset matrix.
    \item The M-step (means calculation) reads the $N \times D$ dataset matrix again along with the $N \times K$ responsibilities.
    \item The M-step (covariance calculation) reads the $N \times D$ dataset matrix a third time to calculate the outer products.
\end{enumerate}

This data-access paradigm forces the processor to stream the massive $N \times D$ memory footprint from main DRAM three separate times per iteration. Consequently, parallel optimization strategies must account for memory movement, as compute pipelines often stall waiting for variable fetches rather than arithmetic exhaustion.

\section{Data Structures \& Memory}
To support the EM algorithm, the implementation relies on several primary data structures. To ensure cache locality and avoid pointer chasing, these are allocated as contiguous 1D memory blocks representing 2D or 3D matrices in row-major order:
\begin{itemize}
    \item \textbf{Dataset (\texttt{data}):} An $N \times D$ matrix containing the input features.
    \item \textbf{Means (\texttt{means}):} A $K \times D$ matrix tracking the centroid of each Gaussian component.
    \item \textbf{Covariances (\texttt{covariances}):} A $K \times D \times D$ tensor representing the covariance matrix for each component.
    \item \textbf{Responsibilities (\texttt{responsibilities}):} An $N \times K$ matrix storing the probability of each point belonging to each class.
    \item \textbf{Inverse Covariances:} A pre-calculated $K \times D \times D$ tensor to avoid repeated matrix inversions inside the inner loop.
\end{itemize}

\textbf{Memory Reference Pattern:} The memory access pattern across the dataset is largely sequential. During the E and M-steps, the outermost loop iterates over the dataset $N$, and the inner loops iterate over $K$ and $D$. Because the dataset is laid out such that the $D$ dimensions for a single point are perfectly contiguous, the CPU accesses the memory sequentially, step-by-step. 

This sequential read pattern is highly cache-friendly. Modern CPU hardware prefetchers can easily detect the linear access stride and proactively load the memory lines into L1 and L2 caches ahead of the execution core, efficiently masking latency. However, for extremely large values of $N$, the full dataset capacity easily exceeds the L3 cache, causing unavoidable capacities misses and cyclic cache flushing over the iterations.

\section{Parallel Modifications}
% Detail any modifications made to the algorithm to run in parallel.
% Mention OpenMP (using threads, collapse, local structures) and CUDA (blocks, grid sizes, atomic operations).

\section{Partitioning Strategy}
% Explain how the data and computations were partitioned.
% e.g., Chunking the N samples across OpenMP threads; mapping threads to samples or to dimensions in CUDA.

\section{Optimizations Overview}
% Provide an overview of your optimized codes.
% What specific optimizations were implemented?
% E.g., CPU Vectorization (AVX2/AVX-512), Private Covariance matrices for lock-free updates (OpenMP), Shared Memory usage (CUDA).
% What problems or challenges did you encounter during optimization?

\section{Experiments and Results}
% Present your experiments and results. Use graphs and tables to show scaling and performance across serial, OpenMP, and CUDA implementations.

\subsection{CPU vs GPU Performance Analysis}
% Specifically from CPU_vs_GPU_Performance.md
During benchmarking, it was observed that the CUDA implementation running on a Tesla P100 GPU is slower than the OpenMP implementation running on a 16-thread or 24-thread CPU for the specific configuration of $D=32$ and $K=8$. This outcome can be attributed to the following architectural factors:

\begin{itemize}
    \item \textbf{Memory-Bound Workload:} GMM training is highly memory-bandwidth bound. The GPU streams the entire dataset ($N \times D$) from HBM2 into SMs three times per iteration (E-Step, M-Step Means, M-Step Covariances). Conversely, modern CPUs with large L3 caches and advanced hardware prefetchers mask memory latency efficiently, keeping the data close to compute units.
    \item \textbf{Synchronization and PCIe Overhead:} The EM iterative loop requires constant CPU/GPU synchronization. Moving partial sums, doing Cholesky decomposition on the CPU, and sending data back via the PCIe bus imposes severe latency, compared to a pure CPU process that uses registers and L1 cache across steps.
    \item \textbf{Shared Memory Serialization:} In early CUDA implementations, using \texttt{atomicAdd} by up to 256 threads into the same shared memory locations (e.g., updating \texttt{s\_Sigma}) causes massive serialization and bank conflicts. OpenMP threads, however, allocate private \texttt{local\_cov} arrays and are completely lock-free until the end.
    \item \textbf{CPU Vectorization "Magic Number":} A dimension of $D=32$ is highly optimal for modern CPUs with AVX2 or AVX-512 instructions. Since $32 \times 4 \text{ bytes} = 128 \text{ bytes}$, the compiler can completely unroll the loops using vectorized instructions processing the dimensions perfectly in very few clock cycles.
\end{itemize}

\subsection{Future Work for GPU Acceleration}
To achieve better absolute performance on the GPU that surpasses highly optimized 24-thread CPU performance, the following architectural redesigns would be necessary:
\begin{enumerate}
    \item \textbf{Kernel Fusion:} Combine the E-step and M-step into a single kernel to reduce global memory reads to exactly once per iteration.
    \item \textbf{GPU Cholesky Base:} Offload the Cholesky decomposition completely onto the GPU (using cuSOLVER) to eliminate iteration-blocking \texttt{cudaMemcpy} calls.
    \item \textbf{Warp-Level Reductions:} Replace shared-memory `atomicAdd` calls with warp shuffle instructions (e.g., \texttt{\_\_shfl\_down\_sync}) to perform lock-free parallel reductions within the SMs.
    \item \textbf{Workload Scale Up:} Test with significantly larger parameters such as $K=64$ and $D=128$, which would exceed CPU cache thresholds and showcase the true massive parallel computational strength of the GPU.
\end{enumerate}

\section{Conclusion}
% What worked and what didn't?
% What limits did you hit?
% Summarize the learning outcomes, validation against truth references, and final thoughts.

\end{document}
